% =====================================================================
%  Lezione 4 — Equilibri alternativi e apprendimento nei giochi
%  Materiale: Lesson_4_Non_Cooperative_Game_Theory_II.pdf,
%             Lesson4-NonCooperativeGameTheory_with_Learning_in_Games.ipynb
% =====================================================================
\chapter{Equilibri alternativi e apprendimento}

\parteteoria

La lezione 3 ha lasciato quattro concetti di soluzione, ognuno con un limite:
\begin{center}
\renewcommand{\arraystretch}{1.25}
\begin{tabularx}{\linewidth}{@{}l X@{}}
\toprule
Concetto & Limite \\
\midrule
Eliminazione iterata (IESDS) & non tutti i giochi sono risolvibili per dominanza \\
Equilibrio di Nash puro & non tutti i giochi ne hanno uno \\
Equilibrio di Nash misto & esiste sempre nei giochi finiti (Nash, 1950), ma è difficile da
calcolare oltre i giochi piccoli \\
Minimax & vale per i giochi a somma zero, dove si calcola in modo efficiente \\
\bottomrule
\end{tabularx}
\end{center}

Questa lezione risponde a tre domande:
\begin{enumerate}[label=\arabic*.]
  \item \emph{perché} l'equilibrio di Nash misto esiste sempre;
  \item quanto è \emph{difficile} calcolarlo;
  \item che cosa si può fare in alternativa: \kw{equilibri correlati} e \kw{apprendimento}
        nei giochi ripetuti.
\end{enumerate}

% ---------------------------------------------------------------------
\section{Esistenza dell'equilibrio di Nash misto}

\begin{teorema}[Nash, 1950]
Ogni gioco finito ha un equilibrio di Nash misto.
\end{teorema}

\textbf{Idea della dimostrazione}: \emph{ridurre} l'esistenza di un equilibrio misto
all'esistenza di un \kw{punto fisso} di una funzione.

\subsection{Il teorema del punto fisso di Brouwer}

\begin{teorema}[Brouwer, 1911]
Ogni funzione \kw{continua} $f$ che manda un insieme euclideo \kw{compatto} e \kw{convesso}
in sé stesso ha un punto $x_0$ tale che $f(x_0) = x_0$.
\end{teorema}

Ricordando le definizioni:
\begin{itemize}
  \item un insieme di $\R^d$ è \kw{compatto} se è chiuso e limitato;
  \item un insieme $C$ è \kw{convesso} se per ogni $x, y \in C$ e ogni
        $\epsilon \in [0,1]$ vale $\epsilon x + (1-\epsilon)y \in C$: contiene il segmento tra
        due suoi punti qualsiasi.
\end{itemize}

Ogni ipotesi serve. Le slide mostrano un controesempio per ciascuna, sull'intervallo
(o sull'anello) e con il grafico di $f$ confrontato con la bisettrice $y = x$: un punto fisso
è un punto in cui il grafico incontra la bisettrice.

\begin{center}
\renewcommand{\arraystretch}{1.3}
\begin{tabularx}{\linewidth}{@{}l X@{}}
\toprule
Ipotesi tolta & Funzione senza punto fisso \\
\midrule
Insieme limitato & $f(x) = x + 0{,}1$ su $\R$: una retta parallela alla bisettrice \\
Insieme chiuso & $f(x) = x^2$ sull'aperto $(0,1)$: i punti fissi di $x^2$ sono $0$ e $1$,
 entrambi fuori dall'insieme \\
Insieme convesso & una rotazione dell'anello attorno al centro: il centro, che sarebbe il punto
fisso, non sta nell'insieme \\
Funzione continua & $f(x) = x + \tfrac12$ per $x < \tfrac12$, $f(x) = x - \tfrac12$ per
 $x \ge \tfrac12$ su $[0,1]$: il grafico ``salta'' oltre la bisettrice \\
\bottomrule
\end{tabularx}
\end{center}

\subsection{La costruzione di Nash}

Il piano della dimostrazione:
\begin{enumerate}
  \item mostrare che l'insieme $S = S_1 \times S_2 \times \dots \times S_n$ dei profili di
        strategie miste è compatto e convesso;
  \item costruire una funzione continua $f : S \to S$ con $f(\sigma) = \sigma$ se e solo se
        $\sigma$ è un equilibrio misto;
  \item applicare Brouwer a $f$.
\end{enumerate}

\paragraph{Passo 1: l'insieme dei profili.}
Le strategie miste del giocatore $i$ formano il \kw{simplesso standard}
\[
  S_i = \bigl\{ x \in \R^{|A_i|} \;:\; x \ge \mathbf{0},\ \uno^T x = 1 \bigr\},
\]
che è compatto e convesso. Il prodotto cartesiano di compatti è compatto e il prodotto di
convessi è convesso, quindi anche $S$ lo è.

\begin{attenzione}[``compatto e chiuso'']
Le slide scrivono che $S$ è ``convesso, compatto e chiuso'' nel piano della dimostrazione e
poi solo ``compatto e chiuso'' nel passo 1. In $\R^d$ un compatto è già chiuso; ciò che
serve per Brouwer, e che i due teoremi enunciati subito dopo dimostrano, è
\emph{compatto e convesso}.
\end{attenzione}

\paragraph{Passo 2: la funzione.}
L'idea è una funzione che prende un profilo e ne restituisce uno \emph{diverso}, a meno che
il profilo non sia un equilibrio. La best response sarebbe la scelta naturale, ma non è una
funzione (può avere più valori) e può essere discontinua. Si usa invece il \kw{guadagno}
che $i$ otterrebbe passando dalla sua strategia mista all'azione pura $a_i^j$:
\[
  g\bigl(a_i^j, \sigma\bigr) = \max\Bigl\{0,\ u_i\bigl(a_i^j, \sigma_{-i}\bigr) - u_i(\sigma)\Bigr\},
\]
e si sposta la probabilità verso le azioni che fanno guadagnare:
\[
  f(\sigma) = \sigma', \qquad
  \sigma'_{i,a_j} = \frac{\sigma_{i,a_j} + g(a_i^j, \sigma)}
                         {\sum_{k \in A_i} \bigl(\sigma_{i,a_k} + g(a_i^k, \sigma)\bigr)} .
\]
La funzione è continua (composizione di funzioni continue, con denominatore $\ge 1$), manda
$S$ in sé (il risultato è ancora una distribuzione di probabilità) e ha un punto fisso se e
solo se $\sigma$ è un equilibrio misto. Per Brouwer un punto fisso esiste, quindi esiste un
equilibrio. $\square$

\begin{extra}[il passaggio che le slide non dimostrano]
Che un equilibrio sia un punto fisso è facile: all'equilibrio nessuna azione pura migliora,
tutti i $g$ sono 0 e $f(\sigma) = \sigma$. Il viceversa è il passo delicato. Sia $\sigma$
un punto fisso e sia $G_i = \sum_k g(a_i^k, \sigma)$; allora per ogni azione
\[
  \sigma_{i,a_j} = \frac{\sigma_{i,a_j} + g(a_i^j,\sigma)}{1 + G_i}
  \quad\Longleftrightarrow\quad
  \sigma_{i,a_j}\,G_i = g(a_i^j, \sigma).
\]
L'utilità $u_i(\sigma)$ è la media delle $u_i(a_i^j, \sigma_{-i})$ pesata con $\sigma_i$,
quindi almeno un'azione del supporto ha $u_i(a_i^j,\sigma_{-i}) \le u_i(\sigma)$, cioè
$g(a_i^j,\sigma) = 0$. Per quell'azione $\sigma_{i,a_j} > 0$ e $\sigma_{i,a_j} G_i = 0$,
quindi $G_i = 0$: \emph{nessuna} azione pura fa guadagnare $i$, e $\sigma_i$ è una best
response. Vale per ogni $i$, quindi $\sigma$ è un equilibrio.
\end{extra}

\begin{attenzione}[refusi]
Il nome è \emph{Brouwer}, non ``Brower''; e $f$ manda $S$ \emph{in} sé stesso (\en{into}),
non necessariamente \emph{su} tutto $S$ (\en{onto}): Brouwer non richiede la suriettività.
\end{attenzione}

% ---------------------------------------------------------------------
\section{Complessità del calcolo degli equilibri}

\subsection{Equilibri misti: PPAD-completezza}

\begin{itemize}
  \item La riduzione appena vista mostra che trovare un equilibrio misto non è più difficile
        che trovare un punto fisso.
  \item Una riduzione (molto più elaborata) mostra il contrario: trovare un punto fisso non
        è più difficile che trovare un equilibrio misto.
  \item Insieme: trovare un equilibrio misto è \emph{esattamente} difficile quanto trovare
        un punto fisso. Il problema è \kw{PPAD-completo}.
\end{itemize}

\begin{definizione}[La classe PPAD]
\kw{PPAD} (\en{Polynomial Parity Argument on Directed graphs}) contiene i problemi di punto
fisso. Il suo problema rappresentativo è \kw{End of Line}: è dato implicitamente un grafo
orientato $G = (V,E)$ con $|V| = 2^n$ nodi, in cui ogni nodo ha grado entrante e uscente al
più 1, tramite due circuiti
\[
  S : \{0,1\}^n \to \{0,1\}^n \ (\text{successore}),\qquad
  P : \{0,1\}^n \to \{0,1\}^n \ (\text{predecessore}).
\]
Dato un nodo sorgente $s$, si deve trovare un pozzo $t$ (la ``fine della linea'').
\end{definizione}

\begin{osservazione}[perché è difficile e perché la soluzione esiste]
La soluzione esiste sempre: seguendo i successori da $s$ si arriva prima o poi a un nodo
senza successore, perché i nodi sono finiti e non si possono ripetere. Ma il grafo ha
$2^n$ nodi e la linea può essere lunga in modo esponenziale; i circuiti $S$ e $P$ non
dicono dove finisce. È lo stesso tipo di problema di Nash: si sa che la soluzione c'è,
non si sa come trovarla in fretta.
\end{osservazione}

\paragraph{Il lemma di Sperner.}
È la controparte discreta di Brouwer. Si prende un reticolo triangolato e se ne colorano i
nodi di bordo con tre colori secondo la regola delle slide (rosso sul lato sinistro, giallo
in basso, blu in alto e a destra). Comunque si colorino i nodi interni, esiste un
\kw{triangolo tricromatico}, anzi ne esiste un numero \emph{dispari}. Trovarne uno è
PPAD-completo.

L'argomento che lo dimostra è proprio un ``End of Line'': si considerano i triangoli come
nodi e le ``porte'' rosso-giallo come archi, e si attraversa ogni porta tenendo il rosso a
sinistra. Partendo dalla porta sul bordo si cammina di triangolo in triangolo, e il cammino
può finire solo in un triangolo tricromatico. Le slide riassumono la catena di riduzioni:
\[
  \textsc{EoL} \le_p \textsc{Sperner} \le_p \textsc{Brouwer} \le_p \textsc{Nash}.
\]

\begin{osservazione}[come leggere la catena]
$X \le_p Y$ significa ``$X$ si riduce a $Y$'': chi sa risolvere $Y$ sa risolvere $X$, quindi
$Y$ è almeno difficile quanto $X$. La catena dice che Nash è almeno difficile quanto End of
Line, cioè PPAD-\emph{difficile}. La costruzione di Nash (Brouwer applicato a $f$) dà la
direzione opposta, l'appartenenza a PPAD. Insieme: PPAD-completo.
\end{osservazione}

\subsection{Equilibri puri: NP-difficoltà}

\kw{Trovare un equilibrio di Nash puro è NP-difficile.} Lo si dimostra riducendo a esso
il problema della \kw{3-colorabilità}: dato un grafo $G = (V,E)$, decidere se si possono
colorare i nodi con tre colori senza che due nodi adiacenti abbiano lo stesso colore.

Si usano i \kw{giochi su grafo}: i giocatori sono i nodi e il payoff di ciascuno dipende
solo dalle interazioni locali (la sua azione e quelle dei vicini).

\paragraph{Primo tentativo.}
Nodi = giocatori, azioni = $\{R, G, B\}$, payoff $+10$ se il colore scelto non è in
conflitto con nessun vicino, $-10$ se lo è. Se $G$ è 3-colorabile c'è un equilibrio con
benessere massimo; se non lo è c'è comunque un equilibrio, solo con benessere più basso.
Questo dimostra che trovare l'equilibrio \emph{di benessere massimo} è NP-difficile, ma non
che lo sia trovarne \emph{uno}.

\paragraph{La costruzione giusta.}
Si aggiunge un \kw{nodo sfidante} collegato a tutti gli altri:
\begin{itemize}
  \item azioni dei nodi del grafo: $\{R, G, B, X, Y\}$; azioni dello sfidante: $\{X, Y\}$;
  \item payoff di un nodo: $+10$ se sceglie un colore senza conflitti, $-10$ se il colore è
        in conflitto, $0$ se sceglie $X$/$Y$ \emph{copiando} lo sfidante, $-5$ se sceglie
        $X$/$Y$ senza copiarlo;
  \item payoff dello sfidante: $+10$ se nessun nodo lo copia, $-10$ se qualcuno lo copia.
\end{itemize}
Allora il gioco ha un equilibrio puro \emph{se e solo se} $G$ è 3-colorabile.

\begin{extra}[perché funziona]
\begin{itemize}
  \item \emph{Se $G$ è 3-colorabile}: i nodi giocano una colorazione corretta ($+10$
        ciascuno) e lo sfidante una lettera qualsiasi ($+10$, nessuno lo copia). Nessun
        nodo ha interesse a passare a una lettera ($\le 0 < 10$), lo sfidante è già al
        massimo: equilibrio.
  \item \emph{Se non lo è}, si supponga per assurdo un equilibrio. Un nodo che gioca una
        lettera senza copiare prende $-5$ e preferirebbe copiare ($0$); quindi in
        equilibrio chi gioca una lettera copia lo sfidante. Ma se almeno un nodo lo copia,
        lo sfidante prende $-10$ e, cambiando lettera, nessuno lo copierebbe più ($+10$):
        devierebbe. Quindi nessun nodo gioca lettere: tutti giocano colori. Poiché $G$ non
        è 3-colorabile, qualche nodo è in conflitto ($-10$) e preferirebbe copiare lo
        sfidante ($0$). Contraddizione.
\end{itemize}
Verificato con una ricerca esaustiva su tutti i profili: il triangolo e il ciclo da 4 nodi
(3-colorabili) hanno un equilibrio puro, il grafo completo $K_4$ (non 3-colorabile) no.
\end{extra}

\begin{attenzione}[NP-difficile rispetto a che cosa]
Se un gioco è dato in forma normale \emph{esplicita} (la tabella completa dei payoff),
trovare un equilibrio puro è facile: basta controllare ogni cella, in tempo polinomiale
nella dimensione della tabella. La difficoltà nasce perché nei giochi su grafo la
rappresentazione è \emph{compatta}: con $|V|$ nodi e 5 azioni la tabella avrebbe
$5^{|V|}\cdot 2$ celle, ma il gioco si descrive con poche regole locali. Il risultato
vale rispetto a questa descrizione compatta, ed è quella che conta nei problemi reali.
\end{attenzione}

\subsection{Il quadro complessivo}

\begin{itemize}
  \item Trovare un equilibrio di Nash misto è \kw{PPAD-completo}.
  \item Trovare un equilibrio di Nash puro è \kw{NP-difficile}.
  \item Trovare l'equilibrio che massimizza il benessere sociale è \kw{NP-difficile}.
\end{itemize}
Conseguenza: se calcolare un equilibrio è difficile per un algoritmo centralizzato e
ottimizzato, non ci si può aspettare che un sistema di agenti autonomi che interagiscono
lo raggiunga da solo nei giochi generici. È la motivazione per le nozioni alternative di
equilibrio e per l'apprendimento.

% ---------------------------------------------------------------------
\section{Equilibri correlati}

\subsection{L'intuizione: il semaforo}

Il gioco dell'\kw{attraversamento}: due auto arrivano a un incrocio e ciascuna sceglie se
passare (\emph{Go}) o fermarsi (\emph{Stop}).
\begin{center}
\renewcommand{\arraystretch}{1.35}
\begin{tabular}{r|c|c|}
\multicolumn{1}{r}{} & \multicolumn{1}{c}{Go} & \multicolumn{1}{c}{Stop}\\ \cline{2-3}
Go   & \pay{-10}{-10} & \pay{5}{0} \\ \cline{2-3}
Stop & \pay{0}{5}     & \pay{-1}{-1} \\ \cline{2-3}
\end{tabular}
\end{center}
\begin{itemize}
  \item Due equilibri puri: (Go, Stop) con $u = (5, 0)$ e (Stop, Go) con $u = (0, 5)$.
        Sbilanciati.
  \item Un equilibrio misto: entrambi giocano Go con probabilità $\tfrac38$, con
        $u_1 = u_2 = -\tfrac58$. Equo ma pessimo: ci si scontra con probabilità
        $\tfrac{9}{64}$.
\end{itemize}

Nella realtà il problema si risolve con un \kw{semaforo}: un \emph{dispositivo di
coordinamento} che raccomanda un'azione a ciascun giocatore. Lo si modella come una
distribuzione di probabilità sui profili:
\begin{center}
\renewcommand{\arraystretch}{1.35}
\begin{tabular}{r|c|c|}
\multicolumn{1}{r}{} & \multicolumn{1}{c}{Go} & \multicolumn{1}{c}{Stop}\\ \cline{2-3}
Go   & $0$         & $\tfrac12$ \\ \cline{2-3}
Stop & $\tfrac12$  & $0$ \\ \cline{2-3}
\end{tabular}
\end{center}
Se entrambi seguono la raccomandazione ottengono $u_1 = u_2 = 2{,}5$. E \kw{deviare dalla
raccomandazione non conviene}, assumendo che l'altro la segua: chi riceve ``Stop'' sa che
l'altro ha ricevuto ``Go'' e, passando, prenderebbe $-10$ invece di $0$; chi riceve ``Go'' sa
che l'altro si ferma e, fermandosi, prenderebbe $-1$ invece di $5$.

\subsection{Definizione}

\begin{definizione}[Equilibrio correlato]
Un \kw{equilibrio correlato} (CE) è una distribuzione di probabilità sui profili di azioni
\[
  p : A_1 \times A_2 \times \dots \times A_n \to [0,1]
\]
tale che per ogni giocatore $i$ e ogni coppia di azioni $a_i, a_i' \in A_i$
\[
  \sum_{a_{-i} \in A_{-i}} p(a_i, a_{-i})\, u_i(a_i, a_{-i})
  \ \ge\ \sum_{a_{-i} \in A_{-i}} p(a_i, a_{-i})\, u_i(a_i', a_{-i}).
\]
\end{definizione}
In parole: se il dispositivo raccomanda $a_i$ al giocatore $i$, la probabilità condizionata
degli altri profili è proporzionale a $p(a_i, \cdot)$, e contro quella distribuzione seguire
la raccomandazione è almeno buono quanto passare a qualunque $a_i'$.

\begin{teorema}[Esistenza]
Ogni gioco finito ha un equilibrio correlato.
\end{teorema}

\begin{extra}[perché esiste, e che rapporto c'è con Nash]
Ogni equilibrio di Nash misto $\sigma$ è anche un equilibrio correlato: basta prendere
$p(a) = \prod_j \sigma_j(a_j)$, la distribuzione dei profili quando ognuno randomizza per
conto suo. L'esistenza segue quindi dal teorema di Nash. Il viceversa non vale: il semaforo
è un equilibrio correlato ma non è il prodotto di due strategie indipendenti (ha
probabilità nulla su (Go,Go) e su (Stop,Stop), cosa impossibile per un prodotto con
entrambe le marginali $(\tfrac12,\tfrac12)$). Gli equilibri correlati sono un insieme più
grande, e quindi possono contenere esiti migliori.
\end{extra}

\subsection{Le disuguaglianze nell'esempio}

Per il giocatore 1 (riga) le due raccomandazioni danno:
\begin{align*}
  \text{raccomandato Go:}\quad
  & p(G,G)\cdot(-10) + p(G,S)\cdot 5 \ \ge\ p(G,G)\cdot 0 + p(G,S)\cdot(-1)\\
  \text{raccomandato Stop:}\quad
  & p(S,G)\cdot 0 + p(S,S)\cdot(-1) \ \ge\ p(S,G)\cdot(-10) + p(S,S)\cdot 5
\end{align*}
e per il giocatore 2 (colonna):
\begin{align*}
  \text{raccomandato Go:}\quad
  & p(G,G)\cdot(-10) + p(S,G)\cdot 5 \ \ge\ p(G,G)\cdot 0 + p(S,G)\cdot(-1)\\
  \text{raccomandato Stop:}\quad
  & p(G,S)\cdot 0 + p(S,S)\cdot(-1) \ \ge\ p(G,S)\cdot(-10) + p(S,S)\cdot 5
\end{align*}
Semplificando:
\[
  p(G,S) \ge \tfrac53\,p(G,G), \quad
  p(S,G) \ge \tfrac35\,p(S,S), \quad
  p(S,G) \ge \tfrac53\,p(G,G), \quad
  p(G,S) \ge \tfrac35\,p(S,S).
\]
Il semaforo $p(G,G) = p(S,S) = 0$, $p(G,S) = p(S,G) = \tfrac12$ le soddisfa tutte.

\subsection{Calcolo con la programmazione lineare}

Le condizioni sono \kw{lineari} nelle incognite $p(a)$: un equilibrio correlato si trova
risolvendo un programma lineare, cioè in tempo polinomiale nella dimensione del gioco.
Gli equilibri correlati sono quindi \emph{più facili} da calcolare di quelli di Nash. Si può
scegliere l'obiettivo che si preferisce:
\[
\begin{aligned}
  \max\ & \text{[obiettivo a scelta]}\\
  \text{s.t. }\ & \sum_{a_{-i}} p(a_i, a_{-i})\bigl(u_i(a_i, a_{-i}) - u_i(a_i', a_{-i})\bigr) \ge 0
  && \forall i,\ \forall a_i, a_i' \in A_i\\
  & \textstyle\sum_{a \in A} p(a) = 1,\qquad p(a) \ge 0 .
\end{aligned}
\]
Due obiettivi tipici:
\begin{itemize}
  \item \kw{benessere totale}: $\max \sum_{a \in A} p(a) \sum_i u_i(a)$;
  \item \kw{regola maximin}: $\max \min_i \sum_{a \in A} p(a)\,u_i(a)$, cioè rendere il più
        alto possibile il payoff atteso del giocatore che sta peggio.
\end{itemize}

\begin{attenzione}[tre errori nelle slide del programma lineare]
\begin{itemize}
  \item Nel vincolo, a destra le slide scrivono $p(a_i', a_{-i})\,u_i(a_i', a_{-i})$. La
        probabilità corretta è $p(a_i, a_{-i})$, la stessa del lato sinistro: si sta
        confrontando che cosa succede \emph{con la stessa raccomandazione} se si gioca
        $a_i$ oppure $a_i'$. Manca anche la sommatoria a destra. La definizione data
        qualche slide prima è giusta; lo è anche il notebook del docente.
  \item L'obiettivo ``$\max \sum_i p(a)\,u_i(a)$'' manca della somma sui profili $a$.
  \item ``$\max\min_i p(a)\,u_i(a)$'' non è un programma lineare così com'è (il minimo non
        è lineare). Lo diventa con una variabile ausiliaria $t$:
        $\max t$ s.t.\ $t \le \sum_a p(a)\,u_i(a)$ per ogni $i$.
\end{itemize}
\end{attenzione}

\section{Altri concetti di soluzione}

\begin{itemize}
  \item \kw{Trattare il gioco come a somma zero} e calcolare una strategia minimax: approccio
        conservativo, utile se si crede che l'avversario voglia danneggiarci o se non si
        conosce la sua utilità.
  \item \kw{Minimizzare il massimo rimpianto}. Il \kw{rimpianto} (\en{regret}) di un'azione
        è quanto si perde rispetto alla migliore risposta, a posteriori:
        \[
          R(a_i, a_{-i}) = \Bigl[\max_{a_i' \in A_i} u_i(a_i', a_{-i})\Bigr] - u_i(a_i, a_{-i}),
          \qquad
          a_i^* \in \argmin_{a_i \in A_i}\ \max_{a_{-i} \in A_{-i}} R(a_i, a_{-i}).
        \]
        È meno conservativo del minimax, utile quando l'avversario è imprevedibile più che
        ostile.
\end{itemize}

\begin{extra}[minimax e minimo rimpianto nell'attraversamento]
Per il giocatore 1: contro Go la risposta migliore vale $0$ (Stop), contro Stop vale $5$
(Go). I rimpianti sono:
\begin{center}
\renewcommand{\arraystretch}{1.2}
\begin{tabular}{lccc}
\toprule
 & contro Go & contro Stop & rimpianto massimo \\
\midrule
Go   & $0 - (-10) = 10$ & $5 - 5 = 0$ & $10$ \\
Stop & $0 - 0 = 0$ & $5 - (-1) = 6$ & $\mathbf{6}$ \\
\bottomrule
\end{tabular}
\end{center}
Il minimo rimpianto massimo è Stop. Anche il criterio maximin sceglie Stop (caso peggiore
$-1$ contro $-10$). Qui i due criteri coincidono, ma in generale no: il maximin guarda solo
i payoff assoluti, il rimpianto li confronta con quello che si sarebbe potuto ottenere.
\end{extra}

% ---------------------------------------------------------------------
\section{Apprendimento nei giochi}

Si immagini ora che i giocatori giochino lo stesso gioco $G$ \emph{ripetutamente}, per $T$
turni. Si vogliono algoritmi adattivi con cui i giocatori \kw{imparino} a giocare strategie
di equilibrio. Poiché calcolare un equilibrio di Nash è PPAD-completo, non ci si aspetta di
impararlo in \emph{ogni} gioco. Due algoritmi classici:
\begin{itemize}
  \item \kw{fictitious play} (Robinson, 1951);
  \item \kw{regret matching} (Hart e Mas-Colell, 2000).
\end{itemize}

\subsection{Fictitious play}

Ogni giocatore stima la strategia degli avversari dalle frequenze osservate e gioca la best
response a quella stima.
\begin{enumerate}
  \item A ogni turno aggiorna il conteggio di quante volte gli avversari hanno giocato
        ciascuna azione: $\text{count}_{a_{-i}}[a_j]$.
  \item Normalizza il conteggio in una distribuzione, la stima della strategia avversaria:
        \[
          \sigma_{-i}^t[a_j] = \frac{\text{count}_{a_{-i}}[a_j]}{t - 1}.
        \]
  \item Gioca una best response alla stima:
        $a_i^t \in \argmax_{a_i} \Ev\bigl[u_i(a_i, a_{-i}) \mid \sigma_{-i}^t\bigr]$.
\end{enumerate}

Proprietà:
\begin{itemize}
  \item aggiorna le credenze assumendo che gli altri giocatori siano \kw{stazionari}
        (giochino sempre la stessa strategia mista), il che non è vero: anche loro imparano;
  \item non è garantito che converga; ma \emph{se} le frequenze $\sigma_{-i}^t$ convergono,
        convergono a un equilibrio di Nash;
  \item converge sicuramente a un equilibrio di Nash nei giochi a due giocatori a somma zero,
        nei giochi risolvibili per eliminazione iterata delle strategie strettamente
        dominate e nei giochi a due giocatori $2\times N$.
\end{itemize}

\begin{osservazione}[che cosa converge]
Le \emph{azioni} giocate in fictitious play sono sempre pure e possono ciclare per sempre:
in sasso-carta-forbici l'algoritmo gioca sasso per un po', poi carta, poi forbici, \ldots\ Ciò
che converge all'equilibrio $(\tfrac13,\tfrac13,\tfrac13)$ sono le \emph{frequenze
empiriche}, cioè le medie nel tempo. Il grafico delle slide mostra proprio queste
oscillazioni sempre più strette attorno a $\tfrac13$.
\end{osservazione}

\subsection{Regret matching}

Ogni giocatore decide quanto spesso giocare un'azione in base a \emph{quanto avrebbe
guadagnato in più} giocandola sempre, col senno di poi.
\begin{enumerate}
  \item A ogni turno calcola il \kw{rimpianto cumulato} di ciascuna azione $a$ rispetto a
        ciò che ha effettivamente giocato:
        \[
          R_i^t(a) = \sum_{s \le t} \Bigl( u_i\bigl(a, \sigma_{-i}^s\bigr)
                     - u_i\bigl(\sigma_i^s, \sigma_{-i}^s\bigr) \Bigr).
        \]
  \item Aggiorna la strategia in proporzione alla \emph{parte positiva} del rimpianto:
        \[
          \sigma_i^{t+1}(a) = \frac{\max(0, R_i^t(a))}{\sum_{a'} \max(0, R_i^t(a'))},
        \]
        e gioca in modo uniforme se tutti i rimpianti sono $\le 0$.
\end{enumerate}

\begin{attenzione}[la formula delle slide è un'altra cosa]
Le slide scrivono $R_i^t(a) = \sum_{s \le t} \max\bigl(0,\ u_i(a,\sigma_{-i}^s) -
u_i(\sigma^s)\bigr)$: la somma delle parti positive di ogni turno, invece della parte
positiva della somma. Non è la stessa cosa: con la versione delle slide i rimpianti non
possono mai diminuire, e un'azione che è stata buona in passato e poi pessima continua a
essere giocata. Il regret matching di Hart e Mas-Colell, quello con le garanzie
teoriche, usa $\max(0, \sum_s \dots)$.

L'effetto si misura nel laboratorio. In sasso-carta-forbici entrambe le versioni sembrano
funzionare, perché per simmetria finiscono comunque attorno a $\tfrac13$. Su un gioco a
somma zero asimmetrico, con strategia minimax $(0{,}33;\ 0{,}56;\ 0{,}11)$, la versione
corretta converge a $(0{,}34;\ 0{,}56;\ 0{,}10)$, quella delle slide a
$(0{,}41;\ 0{,}35;\ 0{,}24)$, che non è un equilibrio. Le slide inoltre non dicono che
cosa fare quando tutti i rimpianti sono nulli (denominatore zero): si gioca uniforme.
\end{attenzione}

Proprietà secondo le slide:
\begin{itemize}
  \item nei giochi generali la media nel tempo converge all'insieme degli
        \kw{equilibri correlati};
  \item converge a un equilibrio di Nash nelle stesse classi di giochi di fictitious play
        (somma zero a due giocatori, risolvibili per dominanza, $2\times N$).
\end{itemize}

\begin{attenzione}[che cosa converge, e a che cosa]
Due precisazioni sull'enunciato delle slide, utili se all'esame si chiede ``a che cosa
converge il regret matching''.
\begin{itemize}
  \item \emph{Che cosa converge.} Le slide parlano della media delle strategie,
        $\bar\sigma^T = \tfrac1T\sum_t \sigma_t$. Ma un equilibrio correlato è una
        distribuzione sui \emph{profili congiunti}, non una strategia per giocatore: ciò
        che converge è la \emph{distribuzione empirica dei profili giocati}, cioè la
        frazione di turni in cui si è giocato ciascun $(a_1,\dots,a_n)$.
  \item \emph{A che cosa.} La versione vista sopra confronta ogni azione con la strategia
        giocata nel suo complesso (rimpianto \emph{incondizionato}, o ``esterno''). La
        teoria garantisce allora la convergenza agli \kw{equilibri correlati grossolani}
        (\en{coarse correlated equilibria}): distribuzioni da cui nessun giocatore vorrebbe
        sganciarsi \emph{prima} di ricevere la raccomandazione, per giocare sempre una
        stessa azione fissa. La convergenza agli equilibri correlati veri e propri vale per
        la variante del paper di Hart e Mas-Colell con rimpianti \emph{condizionati}, che
        confronta ogni azione effettivamente giocata con ciascuna alternativa (``avrei
        fatto meglio a giocare $k$ \emph{ogni volta che} ho giocato $j$?'').
\end{itemize}
\end{attenzione}

\begin{extra}[perché nei giochi a somma zero si arriva comunque a Nash]
Nei giochi a due giocatori a somma zero la distinzione non conta: le marginali di qualunque
equilibrio correlato, anche grossolano, sono strategie minimax. Se nessuno dei due ha
rimpianti verso un'azione fissa, ciascuno ottiene almeno il valore del gioco; poiché la somma
dei payoff è zero, entrambi ottengono esattamente il valore, e le frequenze di ciascuno sono
una strategia ottima. Analogamente, un giocatore senza rimpianti gioca sempre più di rado
un'azione strettamente dominata da un'altra azione pura (ogni volta che la gioca accumula
rimpianto verso quella che la domina), il che spiega il caso dei giochi risolvibili per
dominanza.
\end{extra}

% =====================================================================
\partelab{Equilibri correlati e apprendimento}
% =====================================================================
\fonte{Lesson4-NonCooperativeGameTheory\_with\_Learning\_in\_Games.ipynb}

Il notebook ripete la prima parte del laboratorio 3 (best response, equilibri, minimax,
equilibri $2\times2$: vedi lezione 3) e aggiunge tre esercizi: il programma lineare per gli
equilibri correlati, fictitious play e regret matching. Il notebook contiene gli output
della soluzione del docente: tutti i numeri sotto li riproducono.

\subsection{Equilibri correlati con \texttt{linprog}}

Le incognite sono le $m\cdot n$ probabilità $p(i,j)$, appiattite riga per riga. Ogni coppia
(azione raccomandata, deviazione) dà un vincolo. \texttt{linprog} vuole vincoli $\le 0$ e un
problema di minimo, quindi si cambia segno.

\begin{lstlisting}
import numpy as np
from scipy.optimize import linprog

def build_ce_lp(A1, A2):
    A1, A2 = np.asarray(A1, float), np.asarray(A2, float)
    m, n = A1.shape
    idx = lambda i, j: i * n + j
    rows = []
    for i in range(m):                    # giocatore 1: raccomandato i, devia a k
        for k in range(m):
            if k != i:
                r = np.zeros(m * n)
                for j in range(n):
                    r[idx(i, j)] = -(A1[i, j] - A1[k, j])
                rows.append(r)
    for j in range(n):                    # giocatore 2: raccomandato j, devia a l
        for l in range(n):
            if l != j:
                r = np.zeros(m * n)
                for i in range(m):
                    r[idx(i, j)] = -(A2[i, j] - A2[i, l])
                rows.append(r)
    A_ub, b_ub = np.array(rows), np.zeros(len(rows))
    A_eq, b_eq = np.ones((1, m * n)), np.array([1.0])
    c = -(A1 + A2).flatten()              # max benessere = min -benessere
    bounds = [(0, None)] * (m * n)
    return c, A_ub, b_ub, A_eq, b_eq, bounds, (m, n)

def solve_ce_max_social_welfare(A1, A2, method='highs'):
    c, A_ub, b_ub, A_eq, b_eq, bounds, (m, n) = build_ce_lp(A1, A2)
    res = linprog(c, A_ub=A_ub, b_ub=b_ub, A_eq=A_eq, b_eq=b_eq,
                  bounds=bounds, method=method)
    dist = res.x.reshape(m, n)
    return {"distribution": dist,
            "social_welfare": float((dist * (A1 + A2)).sum()),
            "linprog_result": res}
\end{lstlisting}

\paragraph{Risultati.}
\begin{itemize}
  \item \emph{Attraversamento}: benessere massimo $5{,}0$, distribuzione
        $\bigl(\begin{smallmatrix}0&1\\0&0\end{smallmatrix}\bigr)$, cioè l'equilibrio puro
        (Go, Stop). La matrice dei vincoli è identica a quella del notebook:
        \[
          A_{ub} = \begin{pmatrix}
            10 & -6 & 0 & 0\\ 0 & 0 & -10 & 6\\ 10 & 0 & -6 & 0\\ 0 & -10 & 0 & 6
          \end{pmatrix}
          \qquad\text{(variabili: } p_{GG}, p_{GS}, p_{SG}, p_{SS}\text{)}.
        \]
  \item \emph{Rigori}: benessere $0$, distribuzione uniforme $0{,}25$ su ogni cella.
\end{itemize}

\begin{osservazione}[il programma lineare sceglie un vertice]
Anche il semaforo ha benessere $5$ (verificato: è un equilibrio correlato), ma il risolutore
restituisce (Go, Stop): quando più soluzioni hanno lo stesso valore ottimo, ne restituisce
una qualsiasi, di solito un vertice del poliedro. Se si vuole la soluzione \emph{equa}
bisogna chiederla, con l'obiettivo maximin: risolvendo $\max t$ s.t.\
$t \le \Ev_p[u_i]$ si ottiene esattamente il semaforo,
$\bigl(\begin{smallmatrix}0&\frac12\\\frac12&0\end{smallmatrix}\bigr)$, con $2{,}5$ a testa.
\end{osservazione}

\begin{osservazione}[il commento sui rigori]
Il notebook motiva la distribuzione uniforme con ``il suo equilibrio di Nash è
$[0{,}5; 0{,}5]$''. Il motivo vero è più netto: nei rigori i quattro vincoli diventano
$p_{12} \le p_{11} \le p_{21} \le p_{22} \le p_{12}$, una catena chiusa che forza tutte le
probabilità a essere uguali. L'equilibrio correlato è \emph{unico}, e l'obiettivo non conta
(in un gioco a somma zero il benessere è $0$ per qualunque distribuzione).
\end{osservazione}

\subsection{Fictitious play}

\begin{lstlisting}
def fictitious_play(P, iterations=1000):
    p1_counts, p2_counts = np.ones(3), np.ones(3)   # 1 per evitare divisioni per 0
    sigma_p1, sigma_p2 = [], []
    for _ in range(iterations):
        p1_prob = p1_counts / p1_counts.sum()        # frequenze empiriche
        p2_prob = p2_counts / p2_counts.sum()
        a1 = np.argmax(P @ p2_prob)                  # BR alla stima sull'altro
        a2 = np.argmax(-P.T @ p1_prob)               # somma zero: payoff -P
        p1_counts[a1] += 1
        p2_counts[a2] += 1
        sigma_p1.append(p1_counts / p1_counts.sum())
        sigma_p2.append(p2_counts / p2_counts.sum())
    return p1_prob, p2_prob, np.array(sigma_p1), np.array(sigma_p2)
\end{lstlisting}

In sasso-carta-forbici le frequenze si avvicinano a $\tfrac13$, con oscillazioni che si
stringono:
\begin{center}
\renewcommand{\arraystretch}{1.2}
\begin{tabular}{rcc}
\toprule
Turni & Frequenze del giocatore 1 & Scarto massimo da $\tfrac13$ \\
\midrule
100    & $(0{,}359;\ 0{,}301;\ 0{,}340)$ & $0{,}032$ \\
1000   & $(0{,}347;\ 0{,}319;\ 0{,}334)$ & $0{,}014$ \\
10000  & $(0{,}332;\ 0{,}336;\ 0{,}333)$ & $0{,}002$ \\
\bottomrule
\end{tabular}
\end{center}

\begin{osservazione}[\texttt{argmax} e i pareggi]
\texttt{np.argmax} restituisce sempre il \emph{primo} massimo: è una regola di pareggio
deterministica. In sasso-carta-forbici la matrice è antisimmetrica ($-P^T = P$), quindi i due
giocatori calcolano la stessa best response a partire dagli stessi conteggi: giocano sempre
la stessa azione e hanno frequenze identiche. Non è un errore, ma è un artefatto da tenere presente se si usa fictitious play
per confronti nel progetto.
\end{osservazione}

\subsection{Regret matching}

\begin{lstlisting}
def regret_matching(P, T=10000, seed=0):
    rng = np.random.default_rng(seed)
    R1, R2 = np.zeros(3), np.zeros(3)          # rimpianti cumulati
    p1 = p2 = np.ones(3) / 3
    counts1 = np.zeros(3)
    for t in range(T):
        a1, a2 = rng.choice(3, p=p1), rng.choice(3, p=p2)
        counts1[a1] += 1
        R1 += P[:, a2] - P[a1, a2]             # quanto avrei preso con ogni azione
        R2 += -P[a1, :] + P[a1, a2]            # payoff della colonna = -P
        pos1, pos2 = np.maximum(R1, 0), np.maximum(R2, 0)
        p1 = pos1 / pos1.sum() if pos1.sum() > 0 else np.ones(3) / 3
        p2 = pos2 / pos2.sum() if pos2.sum() > 0 else np.ones(3) / 3
    return counts1 / T
\end{lstlisting}

Qui le azioni sono campionate e il rimpianto si calcola sull'azione effettivamente giocata
dall'avversario; le slide usano le strategie miste (utilità attese). Il rimpianto campionato
è una stima non distorta di quello atteso, quindi le due versioni si comportano allo stesso
modo in media.

\paragraph{Risultati.}
In sasso-carta-forbici, dopo 10\,000 turni, la frequenza del giocatore 1 è
$(0{,}335;\ 0{,}332;\ 0{,}333)$, con scarto massimo da $\tfrac13$ tra $0{,}002$ e $0{,}010$
su 5 semi. La versione con la formula delle slide dà scarti simili su questo gioco, ed è per
questo che il difetto non si vede nel notebook. Il confronto su un gioco asimmetrico
\[
  M = \begin{pmatrix} 2 & -1 & 1 \\ -1 & 1 & 0 \\ 1 & 0 & -1 \end{pmatrix},
  \qquad \text{minimax della riga } (0{,}333;\ 0{,}556;\ 0{,}111),\ \text{valore } 0{,}222,
\]
dopo 20\,000 turni su 3 semi:
\begin{center}
\renewcommand{\arraystretch}{1.2}
\begin{tabular}{lc}
\toprule
Versione & Frequenze della riga \\
\midrule
$\max(0, \sum_s r_s)$ (Hart, Mas-Colell) & $(0{,}34;\ 0{,}56;\ 0{,}10)$, $(0{,}34;\ 0{,}56;\ 0{,}11)$, $(0{,}35;\ 0{,}57;\ 0{,}08)$ \\
$\sum_s \max(0, r_s)$ (slide) & $(0{,}41;\ 0{,}35;\ 0{,}24)$, $(0{,}41;\ 0{,}35;\ 0{,}24)$, $(0{,}41;\ 0{,}35;\ 0{,}24)$ \\
\bottomrule
\end{tabular}
\end{center}
La seconda converge in modo stabile, ma al punto sbagliato: contro la seconda colonna la
strategia $(0{,}41;\ 0{,}35;\ 0{,}24)$ ottiene $-0{,}06$, mentre il minimax garantisce
$0{,}222$ contro qualunque colonna.

% ---------------------------------------------------------------------
\begin{riepilogo}
\begin{itemize}
  \item \textbf{Nash (1950)} via \textbf{Brouwer}: l'insieme dei profili misti è compatto
        e convesso, la funzione dei ``guadagni'' è continua e i suoi punti fissi sono
        esattamente gli equilibri.
  \item Trovare un equilibrio misto è \textbf{PPAD-completo} (End of Line, Sperner);
        trovare un equilibrio puro, o quello di benessere massimo, è \textbf{NP-difficile}
        (riduzione da 3-colorabilità con un nodo sfidante).
  \item \textbf{Equilibrio correlato}: una distribuzione sui profili da cui a nessuno
        conviene deviare data la propria raccomandazione. Ogni Nash è un CE; il semaforo è
        un CE che non è Nash. Si calcola con un programma lineare.
  \item Altri criteri: minimax (avversario ostile), minimo rimpianto massimo
        (avversario imprevedibile).
  \item \textbf{Fictitious play}: best response alle frequenze osservate; converge nei
        giochi a somma zero, risolvibili per dominanza, $2\times N$.
        \textbf{Regret matching}: si gioca in proporzione al rimpianto cumulato positivo;
        la distribuzione dei profili converge agli equilibri correlati (grossolani, per la
        versione con rimpianti incondizionati), e al minimax nei giochi a somma zero.
  \item \textbf{Lab}: CE con \texttt{linprog} (semaforo con l'obiettivo maximin),
        fictitious play e regret matching su RPS; la formula del rimpianto delle slide
        converge al punto sbagliato su un gioco asimmetrico.
\end{itemize}
\textbf{Prossima lezione}: giochi coalizionali (core, nucleolo, valore di Shapley).
\end{riepilogo}
